\section{Concluding remarks}\label{sec:final}
\noindent\textbf{Bounded colourings of $K_n$.}
Theorem~\ref{thm:bounded} leaves open how many trees can be guaranteed exactly.

\begin{question}
What is the largest $f(n)$ such that every globally $\lfloor n/2\rfloor$-bounded edge-colouring of $K_n$ contains $f(n)$ edge-disjoint rainbow spanning trees? Is $\lfloor n/2\rfloor-f(n)$ bounded?
\end{question}

The error term $O(\log^2n)$ comes from our method: the forests we remove have maximum degree $\Theta(\log n)$, and the matrix martingale stays concentrated only while $\Theta(\log n)$ times this many forests remain. We do not know of a colouring that needs any error term at all.

\smallskip

\noindent\textbf{Proper colourings.}
By Theorem~\ref{thm:bounded}, proving Conjecture~\ref{conj:KKS} for large $n$ amounts to finding the last $O(\log^2 n)$ trees, and this must use properness.
Theorem~\ref{thm:gap} and the result of~\cite{GKMO} do this when every colour class has at most $n/2-Cn^{2/3}(\log n)^{2/3}$ edges, and when every colour class has exactly $n/2$ edges, respectively.
Neither covers colourings in which many colour classes have between $n/2-O(n^{2/3}(\log n)^{2/3})$ and $n/2-1$ edges and together form a dense graph, and it would be interesting to prove Conjecture~\ref{conj:KKS} for these.
The exponent $2/3$ in Theorem~\ref{thm:gap} comes from balancing two requirements on the size of the blocks (see Remark~\ref{rem:twothirds}), and we do not know whether it is a limit of the method.

\smallskip

\noindent\textbf{Rainbow forests.}
In Theorem~\ref{thm:forests}, the loss of $\eps m$ forests comes only from the case where $G$ has no orientation gap, which we handle following Xu~\cite{Xu}.
For $K_n$, Theorem~\ref{thm:bounded} shows that the loss can be much smaller, and it would be interesting to know whether every globally $m$-bounded graph with $km$ edges, $k\le2m$, contains $m-O(\log^2m)$ edge-disjoint rainbow $k$-forests.
